% title:          Probability of a square sum
% area:           discrete-probability, enumeration, asymptotics
% methods:        squeeze
% difficulty:
% difficulty-ai:  medium
% status:         partial
% review:         none
% --- history: all optional, leave EMPTY if unknown ---
% origin:         putnam
% origin-date:    1982
% origin-number:  B3
% also-in:
% taken-from:
% references:     43rd Putnam 1982, Problem B3 (answer 4(sqrt2 - 1)/3) - https://prase.cz/kalva/putnam/putn82.html ; Kalva solution - https://prase.cz/kalva/putnam/psoln/psol829.html ; PutnamBench formalisation putnam_1982_b3 (official-style wording, limit in the form r(sqrt s - t)) - https://raw.githubusercontent.com/trishullab/PutnamBench/main/lean4/src/putnam_1982_b3.lean
% history:        ai
\begin{problem}
Let $n\geq 1$ be an integer. Assume $C$ and $D$ are chosen at random from $\{1,\ldots,n\}$. Let $p_n$ be the probability that $C + D$ is a perfect square. Compute $\lim_{n \to \infty}(\sqrt{n} \cdot p_n)$. Express the result in the form $(a\sqrt{b} + c)/d$, where $a$, $b$, $c$, $d$ are integers.
\end{problem}
\begin{solution}
Since for each $n\in\mathbb{N}^{*}$ the number of ways of writing n as the sum of two $a,b\in\mathbb{N}^{*}$ is $n-1$ (indeed fix $a$ it must be between $1$ and $n-1$ for $b$ to be determined so there are as many soltuions as there are as many $a$ i.e. $n-1$). Hence the number of ways of writing as a sum of two positive integer every perfect square between 1 and 2n is:
$$S:=\left(\sum_{i=1}^{\lfloor\sqrt{n}\rfloor} (i^2-1)\right) + \left(\sum_{i=\lfloor\sqrt{n}\rfloor+1}^{\lfloor\sqrt{2n}\rfloor} (n+1) -(i^2-n)\right)$$
$$=\left(\frac{\lfloor\sqrt{n}\rfloor(\lfloor\sqrt{n}\rfloor+1)(2\lfloor\sqrt{n}\rfloor+1)}{6} - \lfloor\sqrt{n}\rfloor\right) + \left(\sum_{i=\lfloor\sqrt{n}\rfloor+1}^{\lfloor\sqrt{2n}\rfloor} (2n+1-i^2)\right) $$
$$=\left(\frac{\lfloor\sqrt{n}\rfloor(\lfloor\sqrt{n}\rfloor+1)(2\lfloor\sqrt{n}\rfloor+1)-6\lfloor\sqrt{n}\rfloor + 6(2n+1)\left(\lfloor\sqrt{2n}\rfloor-\lfloor\sqrt{n}\rfloor-1\right)}{6}\right)-\sum_{i=\lfloor\sqrt{n}\rfloor + 1}^{\lfloor\sqrt{2n}\rfloor}i^2$$
$$=\left(\frac{\lfloor\sqrt{n}\rfloor(\lfloor\sqrt{n}\rfloor+1)(2\lfloor\sqrt{n}\rfloor+1)-6\lfloor\sqrt{n}\rfloor + 6(2n+1)\left(\lfloor\sqrt{2n}\rfloor-\lfloor\sqrt{n}\rfloor-1\right)}{6}\right)$$
$$- \left(\frac{\lfloor\sqrt{2n}\rfloor(\lfloor\sqrt{2n}\rfloor+1)(2\lfloor\sqrt{2n}\rfloor+1)}{6}\right)+\left(\frac{\lfloor\sqrt{n}\rfloor(\lfloor\sqrt{n}\rfloor+1)(2\lfloor\sqrt{n}\rfloor+1)}{6}\right)$$
Then:
$$\sqrt{n}p_n = \sqrt{n}\frac{S}{n^2}=\frac{S}{n^{\frac{3}{2}}}$$
Taking limit and using the classic bounds: $$\sqrt{n}-1<\lfloor\sqrt{n}\rfloor\leq \sqrt{n}$$ we have that the first term tends to, the second to, the third one to. So we sum and obtain: 
\end{solution}
